\documentclass{article}

% -----------------------------
% Packages
% -----------------------------
\usepackage{fancyhdr}
\usepackage{extramarks}
\usepackage{amsmath}
\usepackage{amsthm}
\usepackage{amsfonts}
\usepackage{tikz}
\usepackage[plain]{algorithm}
\usepackage{algpseudocode}
\usepackage{amssymb}

\usetikzlibrary{automata,positioning}

% -----------------------------
% Basic Document Settings
% -----------------------------
\topmargin=-0.45in
\evensidemargin=0in
\oddsidemargin=0in
\textwidth=6.5in
\textheight=9.0in
\headsep=0.25in

\linespread{1.1}

\pagestyle{fancy}
\lhead{\hmwkAuthorName}
\chead{\hmwkClass: \hmwkTitle}
\rhead{}
\lfoot{\lastxmark}
\cfoot{\thepage}

\renewcommand\headrulewidth{0.4pt}
\renewcommand\footrulewidth{0.4pt}

\setlength\parindent{0pt}

% -----------------------------
% Problem Section Helpers
% -----------------------------
\newcommand{\enterProblemHeader}[1]{
  \nobreak\extramarks{}{Problem #1 continued on next page\ldots}\nobreak{}
  \nobreak\extramarks{Problem #1 (continued)}{Problem #1 continued on next page\ldots}\nobreak{}
}

\newcommand{\exitProblemHeader}[1]{
  \nobreak\extramarks{Problem #1 (continued)}{Problem #1 continued on next page\ldots}\nobreak{}
  \nobreak\extramarks{Problem #1}{}\nobreak{}
}

\setcounter{secnumdepth}{0}
\newcounter{partCounter}
\newcounter{homeworkProblemCounter}
\setcounter{homeworkProblemCounter}{1}
\nobreak\extramarks{Problem \arabic{homeworkProblemCounter}}{}\nobreak{}

% -----------------------------
% Theorem env (optional)
% -----------------------------
\newtheorem*{theorem}{Theorem}

% -----------------------------
% Homework Problem Environment
%   Usage:
%   \begin{homeworkProblem}            % auto-number
%     % Your answer here
%   \end{homeworkProblem}
%
%   \begin{homeworkProblem}[1.2.3]     % label problem "1.2.3" (e.g. textbook numbering)
%     % Your answer here
%   \end{homeworkProblem}
% -----------------------------
\newenvironment{homeworkProblem}[1][]{
  \def\hwProblemLabel{#1}
  \ifx\hwProblemLabel\empty
    \edef\hwProblemLabel{\arabic{homeworkProblemCounter}}
    \stepcounter{homeworkProblemCounter}
  \fi
  \section{Problem \hwProblemLabel}
  \setcounter{partCounter}{1}
  \enterProblemHeader{\hwProblemLabel}
}{
  \exitProblemHeader{\hwProblemLabel}
}

% -----------------------------
% Homework Details (EDIT THESE)
% -----------------------------
\newcommand{\hmwkTitle}{Homework \#2}
\newcommand{\hmwkClass}{Math 3320}
\newcommand{\hmwkAuthorName}{\textbf{Joseph Quinn}}

% -----------------------------
% Title
% -----------------------------
\title{
  \vspace{2in}
  \textmd{\textbf{\hmwkClass:\ \hmwkTitle}}\\
  \vspace{3in}
}
\author{\hmwkAuthorName}
\date{}

% -----------------------------
% Convenience Macros (optional)
% -----------------------------
\renewcommand{\part}[1]{\textbf{\large Part \Alph{partCounter}}\stepcounter{partCounter}\\}

% Algorithms
\newcommand{\alg}[1]{\textsc{\bfseries \footnotesize #1}}

% Calculus
\newcommand{\deriv}[1]{\frac{\mathrm{d}}{\mathrm{d}x} (#1)}
\newcommand{\pderiv}[2]{\frac{\partial}{\partial #1} (#2)}
\newcommand{\dx}{\mathrm{d}x}

% Statistics
\newcommand{\E}{\mathrm{E}}
\newcommand{\Var}{\mathrm{Var}}
\newcommand{\Cov}{\mathrm{Cov}}
\newcommand{\Bias}{\mathrm{Bias}}

% Proof step (for aligned derivations)
\newcommand{\step}[2]{& #1 & & \text{#2} \\}

% Solution header (optional)
\newcommand{\solution}{\textbf{\large Solution}}

% -----------------------------
% Document
% -----------------------------
\begin{document}

\maketitle
\pagebreak

\begin{homeworkProblem}[A5]
	\begin{enumerate}
		\item[(I1)]
		      \begin{align*}
			      \text{max}  \quad & 5x_{1} + 6x_{2} + 9x_{3} +8x_{4} \\
			      \text{s.t} \quad  & x_{1}+2x_{2}+3x_{3}+x_{4} \le 5  \\
			                        & x_{1}+x_{2}+2x_{3}+3x_{4} \le 3  \\
			                        & x_{1},x_{2},x_{3},x_{4} \ge 0
		      \end{align*}
		\item[(I2)]
		      \begin{align*}
			      x_{5} & = 5 -x_{1}-2x_{2}-3x_{3}-x_{4}     \\
			      x_{6} & = 3 -x_{1}-x_{2}-2x_{3}-3x_{4}     \\
			      z     & = 5x_{1} + 6x_{2} + 9x_{3} +8x_{4}
		      \end{align*}
		\item[(I3)] The BFS for this dictionary is $x_{1}=0,x_{2}=0,x_{3}=0,x_{4}=0,x_{5}=5,x_{6}=3$ resulting in $z=0$
		\item[(P1)] All of our coefficients are positive meaning we can select any to be our entering variable, using Bland's rule results in $x_{1}$ being our entering variable.
		\item[(P2)] Calculating the ratios for the leaving variable results in:
		      \begin{align*}
			      x_{1} \le 5 \quad \text{for } x_{5} \\
			      x_{1} \le 3 \quad \text{for } x_{6}
		      \end{align*}
		      Meaning $x_{6}$ has the tighter constraint and is our leaving variable.
		\item[(P3)] Creating a new dictionary moving $x_{1}$ to a basic variable and $x_{6}$ to a non basic variable results in the new dictionary,
		      \begin{align*}
			      x_{1} & = 3-x_{2}-2x_{3}-3x_{4}-x_{6}  \\
			      x_{5} & = 2-x_{2}-x_{3}+2x_{4}+x_{6}   \\
			      z     & = 15+x_{2}-x_{3}-7x_{4}-5x_{6}
		      \end{align*}
		\item[(P4)] The BFS for this dictionary is $x_{2}=0, x_{3}=0, x_{4}=0, x_{6}=0, x_{1}=3, x_{5}=2$ resulting in $z=15$


		\item[(P1)] the only positive coefficient in our $z$ euqation is $x_{2}$ so that is our entering vairable
		\item[(P2)] Calculating the ratios for the leaving variable results in:
		      \begin{align*}
			      x_{2} \le 3 \quad \text{for } x_{1} \\
			      x_{2} \le 2 \quad \text{for } x_{5}
		      \end{align*}
		      Meaning $x_{5}$ has the tighter constraint and is our leaving variable.
		\item[(P3)] Creating a new dictionary moving $x_{2}$ to a basic variable and $x_{5}$ to a non basic variable results in the new dictionary,
		      \begin{align*}
			      x_{1} & = 1-x_{3}-5x_{4}+x_{5}-2x_{6}   \\
			      x_{2} & = 2-x_{3}+2x_{4}-x_{5}+x_{6}    \\
			      z     & = 17-2x_{3}-5x_{4}-x_{5}-4x_{6}
		      \end{align*}
		\item[(P4)] The BFS for this dictionary is $x_{3}=0, x_{4}=0, x_{5}=0, x_{6}=0, x_{1}=1, x_{2}=2$ resulting in $z=17$
	\end{enumerate}
	Because our BFS for this recent dicitonary has no positive coefficients in hte $z$ equation this means we cannot pick an entering vairable, also meaning we have arrived at our optimal solution being: \\
	\[
		x_{1}=1, x_{2}=2, x_{3}=0, x_{4}=0
	\]
\end{homeworkProblem}
\begin{homeworkProblem}[A6]
	\begin{enumerate}
		\item[(I1)]
		      \begin{align*}
			      \text{max}  \quad & 3x_{1}+7x_{2}+5x_{3}+x_{4}        \\
			      \text{s.t} \quad  & x_{1}+2x_{2}-4x_{3}+x_{4} \le 7   \\
			                        & 5x_{1}+3x_{2}+x_{3}-8x_{4} \le 10 \\
			                        & x_{1},x_{2},x_{3},x_{4} \ge 0
		      \end{align*}

		\item[(I2)]
		      \begin{align*}
			      x_{5} & = 7- x_{1}-2x_{2}+4x_{3}-x_{4}   \\
			      x_{6} & = 10- 5x_{1}-3x_{2}-x_{3}+8x_{4} \\
			      z     & = 3x_{1}+7x_{2}+5x_{3}+x_{4}
		      \end{align*}
		\item[(I3)] The BFS solution for this dictionary is $x_{1}=0,x_{2}=0,x_{3}=0,x_{4}=0, x_{5}=7, x_{6} = 10$ resulting in $z = 0 $
		\item[(P1)] each coefficient in our $z$ equation is positive so using Bland's rule we select $x_{1}$ to be our entering variable.
		\item[(P2)] calculating the ratios for our basic variable we calculate
		      \begin{align*}
			      x_{1} \le 7  \text{for} \quad x_{5} \\
			      x_{1} \le 2 \text{for} \quad x_{6}
		      \end{align*}
		      Meaning $x_{6}$ has the tigher contraint and becomes our leaving variable.
		\item[(P3)] Creating a new dictionary moving $x_{1}$ to a basic variable and $x_{6}$ to a non basic variable results in the new dictionary,
		      \begin{align*}
			      x_{5} & = 5-\frac{7}{5}x_{2}+\frac{21}{5}x_{3}-\frac{13}{5}x_{4}+\frac{1}{5}x_{6}  \\
			      x_{1} & = 2- \frac{3}{5}x_{2}-\frac{1}{5}x_{3}+\frac{8}{5}x_{4}-\frac{1}{5}x_{6}   \\
			      z     & = 6+\frac{26}{5}x_{2}+\frac{22}{5}x_{3}+\frac{29}{5}x_{4}-\frac{3}{5}x_{6}
		      \end{align*}
		\item[(P4)] Our new dictionary results in a BFS of $x_{1}=2,x_{2}=0,x_{3}=0,x_{4}=0,x_{5}=5,x_{6}=0$ resulting in $z=6$

		\item[(P1)] there are multiple coefficients in our $z$ equation that are positive so using Bland's rule we select $x_{2}$ to be our entering variable.

		\item[(P2)] calculating the ratios for our basic variable we calculate
		      \begin{align*}
			       & x_{2} \le \frac{25}{7}  \text{for} \quad x_{5} \\
			       & x_{2} \le \frac{10}{3}  \text{for} \quad x_{1}
		      \end{align*}
		      Meaning $x_{1}$ has the tighter constraint and becomes our leaving variable.
		\item[(P3)] Creating a new dictionary moving $x_{2}$ to a basic variable and $x_{1}$ to a non basic variable results in the new dictionary,
		      \begin{align*}
			      x_{5} & = \frac{1}{3} + \frac{7}{3}x_1 + \frac{14}{3}x_3 - \frac{19}{3}x_4 + \frac{2}{3}x_6  \\
			      x_{2} & = \frac{10}{3} - \frac{5}{3}x_1 - \frac{1}{3}x_3 + \frac{8}{3}x_4 - \frac{1}{3}x_6   \\
			      z     & = \frac{70}{3} - \frac{26}{3}x_1 + \frac{8}{3}x_3 + \frac{59}{3}x_4 - \frac{7}{3}x_6
		      \end{align*}
		\item[(P4)] This makes our BFS $x_{1}=0, x_{2}=\frac{10}{3}, x_{3}=0, x_{4}= 0, x_{5}=\frac{1}{3}, x_{6}=0$ resulting in $z=\frac{70}{3}$

		\item[(P1)] there are multiple coefficients in our $z$ equation that are positive: $x_{3}$ has coefficient $\frac{8}{3}$ and $x_{4}$ has coefficient $\frac{59}{3}$ ($x_{1}$'s coefficient is negative, so it is not eligible). Using Bland's rule we select $x_{3}$ to be our entering variable.

		\item[(P2)] calculating the ratios for our basic variable we calculate
		      \begin{align*}
			      x_{5} & = \frac{1}{3}+\frac{14}{3}x_{3}: \quad \text{coefficient of } x_{3} \text{ is positive, so } x_{5} \text{ only increases -- no bound} \\
			      x_{2} & = \frac{10}{3}-\frac{1}{3}x_{3} \ge 0 \quad \Longrightarrow \quad x_{3}\le 10
		      \end{align*}
		      Only $x_{2}$'s row gives a binding restriction, so $x_{2}$ becomes our leaving variable at $x_{3}=10$.

		\item[(P3)] Creating a new dictionary moving $x_{3}$ to a basic variable and $x_{2}$ to a non basic variable results in the new dictionary,
		      \begin{align*}
			      x_{3} & = 10-5x_{1}-3x_{2}+8x_{4}-x_{6}     \\
			      x_{5} & = 47-21x_{1}-14x_{2}+31x_{4}-4x_{6} \\
			      z     & = 50-22x_{1}-8x_{2}+41x_{4}-5x_{6}
		      \end{align*}
		\item[(P4)] The BFS for this dictionary is $x_{1}=0,x_{2}=0,x_{4}=0,x_{6}=0,x_{3}=10,x_{5}=47$ resulting in $z=50$

		\item[(P1)] the only positive coefficient in our $z$ equation is $x_{4}$ so it becomes our entering vairable.

		\item[(P2)] calculating the ratios for our basic variable we calculate
		      \begin{align*}
			      x_{3} & = 10+8x_{4}: \quad \text{coefficient of } x_{4} \text{ is positive, so } x_{3} \text{ only increases}  \\
			      x_{5} & = 47+31x_{4}: \quad \text{coefficient of } x_{4} \text{ is positive, so } x_{5} \text{ only increases}
		      \end{align*}
		      Every coefficient of $x_{4}$ in the current dictionary is nonnegative, so there is no ratio bound and thus no leaving variable
	\end{enumerate}
	Since $x_{4}$ can increase without bound while $x_{3}$ and $x_{5}$ remain feasible, the objective $z=50+41x_{4}\to\infty$.  \\
	So this linear program is unbounded.

\end{homeworkProblem}


\end{document}

